\documentclass[12pt]{article}
\usepackage[utf8]{inputenc}
\usepackage[T1]{fontenc}
\usepackage[margin=1in]{geometry}
\usepackage{setspace}
\usepackage{xcolor}
\usepackage[colorlinks=true, urlcolor=blue!60!black, linkcolor=blue!60!black]{hyperref}
\setlength{\parindent}{2em}
\setlength{\parskip}{1.25em}
\renewcommand{\baselinestretch}{1.0}
\begin{document}
\begin{center}
{\LARGE \textit{Hayburn's Case}, 2 U.S. (2 Dall.) 409 (1792)\par}
\vspace{0.8em}
{\large Case Brief --- Four-Page Edition\par}
\end{center}
\vspace{1em}

\section{Memory Jogger}

When Congress told federal circuit courts to decide veterans' pension claims subject to revision by the Secretary of War and Congress, five Justices on circuit refused, writing that Congress may not assign Article III courts non-judicial duties or make their decisions subject to executive or legislative revision; the Supreme Court itself never reached the merits (\href{https://supreme.justia.com/cases/federal/us/2/409/}{Justia}).

\section{Detailed Case Facts}

The Invalid Pensions Act of March 23, 1792, directed disabled Revolutionary War veterans to apply to the federal circuit courts, which were to examine proof of disability, determine eligibility and pension rates, and certify their findings to the Secretary of War. The courts' determinations were not final: the Secretary could withhold a veteran's name whenever he had ``cause to suspect imposition or mistake,'' and Congress retained authority to revise or reject the awards. Because Justices then rode circuit, five of the six sitting Justices confronted the Act at once.

\section{Procedural History}

In April 1792 the Pennsylvania Circuit (Justices Wilson and Blair and Judge Peters) refused to hear veteran William Hayburn's petition and wrote to President Washington that the Act was unconstitutional. The New York and North Carolina circuits wrote similar letters, though the New York judges agreed to hear claims out of court as commissioners. In August 1792 Attorney General Edmund Randolph sought mandamus in the Supreme Court to compel the Pennsylvania court to act. The Court divided 3--3 on his \textit{ex officio} authority and denied that motion; renewed on Hayburn's behalf, the motion was held under advisement until Congress, in February 1793, created a non-judicial pension procedure, mooting the case.

\section{Judicial Votes}

\begin{itemize}
\item \textbf{Supreme Court:} No merits judgment; a 3--3 division denied the \textit{ex officio} motion, and the renewed motion was never decided. The report is a procedural entry by Alexander Dallas, with a footnote reprinting the circuit letters.
\item \textbf{Views on the merits:} Chief Justice Jay and Justices Cushing, Wilson, Blair, and Iredell, in the New York, Pennsylvania, and North Carolina circuit letters.
\item \textbf{Concurrences and dissents:} None in the modern sense.
\end{itemize}

\section{Holding}

No Supreme Court merits holding. But the unanimous view of five Justices in the circuit letters established two enduring Article III principles: Congress may not assign Article III courts non-judicial duties, and federal courts may not render advisory or non-final judgments subject to executive or legislative revision.

\section{Analysis of Opinions}

\textbf{Separation of powers.} All three circuits began from the Constitution's division of the government into three distinct and independent branches, with the judicial power vested exclusively in Article III courts. Neither Congress nor the President may compel courts to perform duties that are not ``properly judicial'' and performed ``in a judicial manner.''

\textbf{Finality.} The core defect was revision. Letting an executive officer set aside courts' findings for suspected imposition or mistake, and letting Congress override them, made courts subordinate advisers to the political branches. The North Carolina judges added that the Legislature holds no judicial power apart from impeachment.

\textbf{Differences among the circuits.} The New York court construed the Act as naming the judges individually as commissioners, free to accept or decline, and sat as commissioners out of respect for Congress and the veterans. The Pennsylvania court refused to act in any capacity. The North Carolina court doubted that judges could constitutionally serve as commissioners at all.

\textbf{Later.} \textit{Plaut v. Spendthrift Farm} (1995) relied on \textit{Hayburn's Case} to invalidate a statute reopening final judgments. \textit{Muskrat v. United States} (1911) applied the advisory-opinion bar, while \textit{Mistretta} (1989) invoked the commissioner theory to uphold judges' service on the Sentencing Commission.

\section{Examples: Future Applications}

\begin{enumerate}
\item \textbf{Same-Side (Invalid): Executive Veto of Awards.} A statute lets the Homeland Security Secretary vacate any district-court cyber-relief award inconsistent with fiscal priorities.
\item \textbf{Same-Side (Invalid): Advisory Rulings on Bills.} The D.C. Circuit must tell committee chairs whether pending bills are constitutional---a duty not properly judicial.
\item \textbf{Opposite-Side (Valid): Ending Prospective Injunctions.} A change in law requiring courts to terminate prospective decrees does not reopen final judgments (\textit{Miller v. French}).
\item \textbf{Opposite-Side (Valid): Voluntary Commission Service.} Judges serving individually on a non-adjudicatory ethics commission follow \textit{Mistretta}.
\item \textbf{Fence-Sitter: Judgments Payable Only by Appropriation.} Veterans' judgments are paid only if Congress funds them, and Congress refuses one. Is that the power of the purse, or legislative suspension of a judgment?
\end{enumerate}

\section{Critique}

\textbf{An advisory opinion against advisory opinions.} The judges protested advisory functions by writing unsolicited constitutional letters to the President outside any lawsuit.

\textbf{The commissioner fiction.} Letting judges adjourn court and decide the same claims in the same room as commissioners invites doubt that formal titles protect independence. And as Maeva Marcus and Robert Teir note, the Court's only ruling was a 3--3 tie; later Courts elevated a reporter's footnote into canon.

\section{Key Quotations}

\begin{quote}
``That by the Constitution of the United States, the government thereof is divided into three distinct and independent branches, and that it is the duty of each to abstain from, and to oppose, encroachments on either. That neither the Legislative nor the Executive branches can constitutionally assign to the Judicial any duties but such as are properly judicial, and to be performed in a judicial manner.''
\end{quote}
\noindent --- New York Circuit Court (Jay, C.J., Cushing, J., Duane, D.J.), 2 U.S. (2 Dall.) at 410 n.*.

\begin{quote}
``Such revision and control we deemed radically inconsistent with the independence of that judicial power which is vested in the courts; and, consequently, with that important principle which is so strictly observed by the Constitution of the United States.''
\end{quote}
\noindent --- Pennsylvania Circuit Court (Wilson, J., Blair, J., Peters, D.J.), 2 U.S. (2 Dall.) at 411 n.*.

\begin{quote}
``And we beg leave to add, with all due deference, that no decision of any court of the United States can, under any circumstances, in our opinion, agreeable to the Constitution, be liable to a reversion, or even suspension, by the Legislature itself, in whom no judicial power of any kind appears to be vested, but the important one relative to impeachments.''
\end{quote}
\noindent --- North Carolina Circuit Court (Iredell, J., Sitgreaves, D.J.), 2 U.S. (2 Dall.) at 413 n.*.

\end{document}
